\documentclass[a4paper,12pt]{article}

% Packages utiles
\usepackage[utf8]{inputenc}
\usepackage[T1]{fontenc}
\usepackage[french]{babel}
\usepackage{graphicx}
\usepackage{geometry}
\usepackage{hyperref}
\usepackage{titlesec}

\geometry{margin=2.5cm}



\begin{document}

\begin{titlepage}
    \thispagestyle{empty}

    \noindent
    \begin{minipage}[t]{0.35\textwidth}
        \vspace{0pt}
        \includegraphics[width=0.9\textwidth]{img/iut_logo.png}
    \end{minipage}
    \hfill
    \begin{minipage}[t]{0.55\textwidth}
        \vspace{0pt}
        \begin{flushright}
            \textbf{Universite Paris Est Creteil} \\
            \textbf{IUT Senart Fontainebleau}
        \end{flushright}
    \end{minipage}

    \vspace{2cm}

    \begin{center}
        {\LARGE \textbf{SAE 4.1  Chuzzle}}\\[0.8cm]
        {\Large Rapport d’avancement}\\[0.8cm]

        \includegraphics[width=0.3\textwidth]{img/logo.png}\\[1cm]

        {\large
        Algassimou Pellel DIALLO \\[0.7cm]
        Ayoub ANHDIRE \\[0.7cm]
        Youness BOULALAM
        }
    \end{center}

    \vfill

    \begin{center}
        \date{\today}
    \end{center}
\end{titlepage}

\tableofcontents
\newpage

%--------------------------------------------------
\section{Introduction}

Dans le cadre de la SAE 4.1, nous avons développé une application Android inspirée du jeu \textbf{Chuzzle}.

Le jeu se base sur une grille de taille 6x6 contenant différents éléments colorés. Le joueur peut effectuer des glissements de lignes ou de colonnes afin de créer des séries d’au moins trois éléments identiques.

Lorsqu’une série est formée :
\begin{itemize}
    \item elle est supprimée,
    \item les éléments au-dessus tombent (gravité),
    \item de nouveaux éléments apparaissent.
\end{itemize}

L’objectif est d’obtenir le meilleur score possible avant qu’aucun coup ne soit possible.

Le projet est structuré selon une architecture \textbf{MVC (Modèle - Vue - Contrôleur)}.

%--------------------------------------------------
\section{Fonctionnalités}

\subsection{Fonctionnalités implémentées}

\begin{itemize}
    \item Génération d’une grille sans série initiale
    \item Détection des séries
    \item Suppression des séries
    \item Application de la gravité
    \item Remplissage des cases vides
    \item Calcul du score
    \item Menu principal
    \item Lancement d’une partie
    \item Affichage de la grille
    \item Détection de fin de partie
\end{itemize}

\subsection{Captures d’écran}

\begin{figure}[h]
    \centering
%  \includegraphics[width=0.5\textwidth]{menu.png}
    \caption{Menu principal}
\end{figure}

% Ajoute tes images ici

%--------------------------------------------------
\section{Structure du programme}

\subsection{Modèle}

\begin{itemize}
    \item Grille
    \item Case
    \item TypeElement
    \item Serie
    \item Position
    \item MoteurJeu
    \item ScoreManager
\end{itemize}

\subsection{Vue}

\begin{itemize}
    \item MenuActivity
    \item JeuActivity
    \item OptionsActivity
    \item CaseView
\end{itemize}

\subsection{Contrôleur}

\begin{itemize}
    \item MenuController
    \item JeuController
\end{itemize}

\subsection{Rendu graphique}

\begin{itemize}
    \item ElementRenderer
    \item NormalRenderer
    \item DaltonienRenderer
\end{itemize}

\subsection{Diagramme UML}

\begin{figure}[h]
    \centering
 %   \includegraphics[width=0.8\textwidth]{uml.png}
    \caption{Diagramme de classes}
\end{figure}

%--------------------------------------------------
\section{Algorithme de fin de partie}

La partie se termine lorsqu’aucun coup n’est possible.

\subsection{Principe}

L’algorithme consiste à tester tous les déplacements possibles sur la grille.

\subsection{Étapes}

\begin{enumerate}
    \item Parcourir toutes les lignes
    \item Tester les décalages possibles
    \item Vérifier si une série est créée
    \item Répéter pour les colonnes
    \item Si aucun coup ne crée de série :
\end{enumerate}

\begin{center}
\textbf{La partie est terminée}
\end{center}

\subsection{Implémentation}

Cette vérification est réalisée dans la méthode :

\begin{verbatim}
Grille.existeCoupPossible()
\end{verbatim}

Puis utilisée dans :

\begin{verbatim}
MoteurJeu.jouerCoup()
\end{verbatim}

%--------------------------------------------------
\section{Conclusion}

\subsection{Algassimou Pellel DIALLO}

Ce projet m’a permis de comprendre :
\begin{itemize}
    \item l’architecture MVC
    \item la logique d’un moteur de jeu
    \item le développement Android
\end{itemize}

Difficultés rencontrées :
\begin{itemize}
    \item gestion des interactions
    \item affichage graphique
\end{itemize}

Améliorations possibles :
\begin{itemize}
    \item animations
    \item optimisation graphique
\end{itemize}

\subsection{Youness BOULALAM}

Ma contribution à ce projet a principalement porté sur la mise en place du générateur
de types aléatoires, le calcul du score de base, l'affichage des informations de jeu
à l'écran ainsi que sur les aspects liés à l'accessibilité visuelle.
La réalisation de l'affichage daltonien a constitué pour moi le défi le plus stimulant
de ce projet. Il ne s'agissait pas uniquement d'un problème technique, mais d'une
réflexion plus large sur la conception d'une interface : différencier les
éléments par la forme autant que par la couleur Suite à vos explications, nous avons constaté que l’utilisation des couleurs seules n’était pas suffisante et qu’il était nécessaire d’ajouter un élément visuel plus distinctif.
Ce projet m'a permis de consolider mes compétences en développement Android,
notamment dans la manipulation du Canvas et la gestion du cycle de vie des activités.
Il m'a également sensibilisé à l'importance d'une architecture bien définie dès le
départ, sans laquelle la coordination entre les membres de l'équipe aurait été
considérablement plus difficile.
À titre d'amélioration, il serait intéressant d'enrichir l'expérience visuelle par
des animations lors de la suppression des séries.
\subsection{Ayoub ANHDIRE}

À compléter.

\section*{Conclusion générale}

Ce projet nous a permis de concevoir et de développer une application Android complète,
de la modélisation du moteur de jeu jusqu'à la réalisation de l'interface utilisateur.
En partant d'une architecture MVC rigoureuse, nous avons progressivement implémenté
l'ensemble des mécaniques du jeu Chuzzle : génération de la grille, détection et
suppression des séries, application de la gravité, gestion des verrous et calcul du score.

Le travail en trinôme, organisé en sprints hebdomadaires, nous a imposé une rigueur
dans la répartition des tâches et dans la communication autour du code. L'utilisation
de Gitea comme outil de versionnement a été essentielle pour assurer la cohérence du
projet tout au long du développement.

Au-delà des fonctionnalités attendues, nous avons porté une attention particulière à
l'accessibilité, notamment pour les joueurs daltoniens, en proposant des formes
distinctes pour chaque type d'élément en complément des couleurs. Le projet constitue
ainsi une expérience complète de développement logiciel, alliant conception,
implémentation et réflexion sur la qualité du produit final.
%--------------------------------------------------

\end{document}
